% Part: many-valued-logic
% Chapter: three-valued-logics
% Section: kleene

\documentclass[../../../include/open-logic-section]{subfiles}

\begin{document}

\olfileid{mvl}{thr}{skl}

\olsection{منطق‌های کلینی}

استیون کلینی دو منطق سه‌ارزشی معرفی کرد که انگیزهٔ آن‌ها منطقی بود که
در آن ارزش‌های صدق به‌منزلهٔ خروجی‌های رویه‌های محاسباتی در نظر گرفته
می‌شوند: یک رویه ممکن است~$\True$ یا~$\False$ به دست دهد، اما همچنین
ممکن است پایان نیابد. در آن صورت ارزش صدق متناظر تعریف‌نشده است و با
$\Undef$ نمایش داده می‌شود.

برای محاسبهٔ نقیض گزارهٔ~$!A$، ابتدا مقدار~$!A$ را محاسبه و سپس نقیض
حاصل را بازمی‌گردانیم. اگر محاسبهٔ~$!A$ پایان نیابد، کل رویه نیز پایان
نمی‌یابد؛ پس نقیض~$\Undef$ همان~$\Undef$ است.

برای محاسبهٔ عطف~$!A \land !B$ دو گزینه وجود دارد. می‌توان ابتدا~$!A$
و سپس~$!B$ را محاسبه کرد؛ حاصل اگر خروجی هر دو~$\True$ باشد~$\True$، و
در غیر این صورت~$\False$ خواهد بود. اگر یکی از دو محاسبه متوقف نشود،
کل رویه نیز متوقف نمی‌شود. بنابراین در این حالت اگر یکی از مؤلفه‌های
عطف تعریف‌نشده باشد، خود عطف نیز تعریف‌نشده است. همین امر دربارهٔ فصل
نیز برقرار است.

اما اگر بتوانیم~$!A$ و~$!B$ را به‌طور موازی ارزیابی کنیم، نتیجهٔ بهتری
می‌گیریم. اگر یکی از دو رویه متوقف شود و~$\False$ بازگرداند، می‌توانیم
دست نگه داریم، زیرا پاسخ باید کاذب باشد. بنابراین در این حالت، عطفی با
یک مؤلفهٔ کاذب کاذب است، حتی اگر مؤلفهٔ دیگر تعریف‌نشده باشد.
به‌همین‌ترتیب، هنگام محاسبهٔ موازی یک فصل، به‌محض اینکه رویهٔ یکی از دو
مؤلفهٔ فصل مقدار صادق را بازگرداند، می‌توانیم متوقف شویم؛ زیرا فصل باید
صادق باشد. پس در این حالت حتی اگر یکی از مؤلفه‌ها تعریف‌نشده باشد،
می‌توانیم حاصل یک ادعای مرکب را بدانیم. با این تفسیر می‌توان~$\Undef$
را به‌جای «تعریف‌نشده»، «ناشناخته» خواند.

این دو تفسیر به منطق قوی و ضعیف کلینی می‌انجامند. رابط شرطی هم‌ارز
$\lnot !A \lor !B$ تعریف می‌شود.

\begin{defn}
\emph{منطق قوی کلینی}~$\LogKs$ با ماتریس زیر تعریف می‌شود:
\begin{enumerate}
  \item زبان گزاره‌ای استاندارد~$\Lang L_0$ با~$\lnot$، $\land$،
  $\lor$ و~$\lif$.
  \item مجموعهٔ ارزش‌های صدق~$V = \{\True, \Undef, \False\}$.
  \item $\True$ تنها ارزش ممتاز است؛ یعنی~$V^+ = \{\True\}$.
  \item توابع صدق با جدول‌های زیر داده می‌شوند:
  \begin{center}
    \begin{tabular}{c|c} 
      $\tf{\lnot}$ & \\ 
      \hline  
      $\True$ & $\False$ \\ 
      $\Undef$ & $\Undef$ \\
      $\False$ & $\True$ 
    \end{tabular}
    \quad
    \begin{tabular}{c|ccc} 
      $\tf{\land}[\LogKs]$ & $\True$ & $\Undef$ & $\False$ \\ 
      \hline 
      $\True$ & $\True$ & $\Undef$ & $\False$ \\ 
      $\Undef$ & $\Undef$ & $\Undef$ & $\False$\\ 
      $\False$ & $\False$ & $\False$ & $\False$ 
    \end{tabular}
    \\[2ex]
    \begin{tabular}{c|ccc} 
      $\tf{\lor}[\LogKs]$ & $\True$ & $\Undef$ & $\False$ \\ 
      \hline 
      $\True$ & $\True$ & $\True$ & $\True$ \\ 
      $\Undef$ & $\True$ & $\Undef$ & $\Undef$ \\
      $\False$ & $\True$ & $\Undef$ & $\False$ 
    \end{tabular}
    \quad
    \begin{tabular}{c|ccc} 
      $\tf{\lif}[\LogKs]$ & $\True$ & $\Undef$ & $\False$ \\ 
      \hline 
      $\True$ & $\True$ & $\Undef$ & $\False$ \\ 
      $\Undef$ & $\True$ & $\Undef$ & $\Undef$  \\ 
      $\False$ & $\True$ & $\True$ & $\True$ 
    \end{tabular}
  \end{center} 
\end{enumerate}
\end{defn}

\begin{defn}
\emph{منطق ضعیف کلینی}~$\LogKw$ با ماتریس زیر تعریف می‌شود:
\begin{enumerate}
  \item زبان گزاره‌ای استاندارد~$\Lang L_0$ با~$\lnot$، $\land$،
  $\lor$ و~$\lif$.
  \item مجموعهٔ ارزش‌های صدق~$V = \{\True, \Undef, \False\}$.
  \item $\True$ تنها ارزش ممتاز است؛ یعنی~$V^+ = \{\True\}$.
  \item توابع صدق با جدول‌های زیر داده می‌شوند:
  \begin{center}
    \begin{tabular}{c|c} 
      $\tf{\lnot}$ & \\ 
      \hline  
      $\True$ & $\False$ \\ 
      $\Undef$ & $\Undef$ \\
      $\False$ & $\True$ 
    \end{tabular}
    \quad
    \begin{tabular}{c|ccc} 
      $\tf{\land}[\LogKw]$ & $\True$ & $\Undef$ & $\False$ \\ 
      \hline 
      $\True$ & $\True$ & $\Undef$ & $\False$ \\ 
      $\Undef$ & $\Undef$ & $\Undef$ & $\Undef$\\ 
      $\False$ & $\False$ & $\Undef$ & $\False$ 
    \end{tabular}
    \\[2ex]
    \begin{tabular}{c|ccc} 
      $\tf{\lor}[\LogKw]$ & $\True$ & $\Undef$ & $\False$ \\ 
      \hline 
      $\True$ & $\True$ & $\Undef$ & $\True$ \\ 
      $\Undef$ & $\Undef$ & $\Undef$ & $\Undef$ \\
      $\False$ & $\True$ & $\Undef$ & $\False$ 
    \end{tabular}
    \quad
    \begin{tabular}{c|ccc} 
      $\tf{\lif}[\LogKw]$ & $\True$ & $\Undef$ & $\False$ \\ 
      \hline 
      $\True$ & $\True$ & $\Undef$ & $\False$ \\ 
      $\Undef$ & $\Undef$ & $\Undef$ & $\Undef$  \\ 
      $\False$ & $\True$ & $\Undef$ & $\True$ 
    \end{tabular}
  \end{center} 
\end{enumerate}
\end{defn}

\begin{prop}
  $\LogKs$ و~$\LogKw$ هیچ همان‌گویی‌ای ندارند.
\end{prop}

\begin{proof}
  اگر برای همهٔ~!!{propositional variable}های~$p$ داشته باشیم
  $\pAssign v(p) = \Undef$، هر فرمول~$!A$ ارزش صدق
  $\pValue v(!A) = \Undef$ خواهد داشت، زیرا در هر دو منطق
  \[
    \tf{\lnot}(\Undef) = \tf{\lor}(\Undef, \Undef) = \tf{\land}(\Undef,
  \Undef) = \tf{\lif}(\Undef, \Undef) = \Undef.
  \]
  چون در هیچ‌یک از~$\LogKs$ و~$\LogKw$، $\Undef \notin V^+$، در این
  !!{valuation}، مقدار~$!A$ ممتاز نخواهد بود.
\end{proof}

هرچند منطق ضعیف و قوی کلینی هیچ همان‌گویی‌ای ندارند، روابط نتیجهٔ
نابدیهی دارند.

\begin{prob}
  کدام‌یک از روابط زیر در (الف) منطق قوی و (ب) منطق ضعیف کلینی
  برقرارند؟ برای هریک جدول صدق بدهید.
  \begin{enumerate}
    \item $p, p \lif q \Entails q$
    \item $p \lor q, \lnot p \Entails q$
    \item $p \land q \Entails p$
    \item $p \Entails p \land p$
    \item $p \Entails p \lor q$
  \end{enumerate}
\end{prob}

دیمیتری بوچوار~$\Undef$ را به معنای «بی‌معنا» تفسیر کرد و کوشید با
مقررکردن اینکه جمله‌های پارادوکسیک مقدار~$\Undef$ بگیرند، از آن برای حل
پارادوکس‌هایی مانند پارادوکس دروغگو استفاده کند. او منطقی معرفی کرد که
در اصل همان منطق ضعیف کلینی است که با رابط‌های دیگری گسترش یافته؛ دو
مورد از آن‌ها «نقیض بیرونی» و عملگر «تعریف‌نشده است» هستند:
\begin{center}
  \begin{tabular}{c|c} 
    $\tf{\sim}$ & \\ 
    \hline  
    $\True$ & $\False$ \\ 
    $\Undef$ & $\True$ \\
    $\False$ & $\True$ 
  \end{tabular}
\quad
  \begin{tabular}{c|c} 
    $\tf{+}$ & \\ 
    \hline  
    $\True$ & $\False$ \\ 
    $\Undef$ & $\True$ \\
    $\False$ & $\False$ 
  \end{tabular}
\end{center}

\begin{prob}
  آیا می‌توانید~$\sim$ را در منطق بوچوار برحسب~$\lnot$ و~$+$ تعریف
  کنید؟ یعنی~!!a{formula} بیابید که فقط~!!{propositional variable}
  چون~$p$ را داشته باشد و شامل~$\sim$ نباشد، اما همواره همان ارزش صدق
  $\mathord{\sim}p$ را بگیرد. برای نشان‌دادن درستی پاسخ جدول صدق بدهید.
\end{prob}


\end{document}
